\section{Pipeline Architecture}
\label{sec:pipeline}

Figure~\ref{fig:pipeline} illustrates the five-stage pipeline.
Each stage is independently configurable and evaluated in the ablation
study of Section~\ref{sec:experiments}.

\begin{figure*}[!tbp]
  \centering
  \includegraphics[width=\linewidth]{fig15_retrieval_overview.pdf}
  \caption[Schema retrieval pipeline overview]{
    The schema retrieval pipeline.
    At query time (top row), the natural-language question is encoded
    by the INT8-quantized arctic-embed-m encoder using an asymmetric
    query prefix; the resulting 768-dimensional vector is matched against
    a pre-computed per-database sqlite-vec BLOB index via cosine
    approximate nearest-neighbour search; the top-$k$ tables and their
    DDL are inserted into the DAIL-SQL prompt for the local LLM.
    At index construction time (bottom row), table DDL is embedded
    once without a prefix and stored; the same INT8 encoder weights are
    reused, so the 91\,MB RAM footprint applies to both paths.
    Retrieval latency (18.4\,ms) represents a 15.2\% reduction over FP32
    with zero quality loss.}
  \label{fig:pipeline}
\end{figure*}

\subsection{Schema Retrieval}

Given a natural-language question $q$ and a target database $d$, the
retriever identifies the $k$ most relevant tables from $d$.
Each table is represented as its DDL \texttt{CREATE TABLE} statement
concatenated with three sample rows.
arctic-embed-m~\citep{merrick2024arctic} embeds the question with an
asymmetric query prefix and each table document without a prefix.
Cosine similarity selects the top-$k$ tables, which are included in the
LLM prompt.

Table embeddings are pre-computed at index construction time and stored in
per-database \texttt{sqlite-vec}~\citep{garcia2024sqlitevec} BLOB indices;
retrieval requires a single matrix--vector product followed by an approximate
nearest-neighbour lookup, with no external vector database server.
The encoder applies dynamic INT8 quantization via
\texttt{torch.quantization.quantize\_dynamic}, targeting all
\texttt{nn.Linear} layers; Section~\ref{sec:quantization} provides a full
analysis of this choice.

\subsection{Few-Shot Example Selection (sql2skeleton)}
\label{sec:fewshot}

Few-shot examples improve SQL generation by providing structural
demonstrations of similar queries.
For each development question, the system retrieves $k$ training questions
from the DAIL-SQL selection pool of 8{,}659 examples~\citep{gao2023dailsql}.

Retrieval uses sql2skeleton similarity: a SQL query is reduced to a
\emph{skeleton} by replacing all string and numeric literals with the
placeholder \texttt{VALUE} and lower-casing function names, while
preserving every structural keyword---\texttt{SELECT}, \texttt{FROM},
\texttt{JOIN}, \texttt{WHERE}, \texttt{GROUP BY}, \texttt{HAVING}, and
nested subqueries.
The four Spider difficulty levels~\citep{yu2018spider} produce qualitatively
different skeletons.
\emph{Easy} queries (\texttt{SELECT}, simple \texttt{WHERE}) collapse to a
one-line skeleton; \emph{Medium} queries preserve \texttt{JOIN} and
\texttt{GROUP BY} structure; \emph{Hard} queries retain multi-table joins and
\texttt{HAVING} clauses; \emph{Extra Hard} queries expose nested
\texttt{NOT IN (SELECT\ldots)} subquery structure.
This selection strategy substantially reduces the occurrence of
\texttt{no\_such\_column} errors relative to question-only similarity
(Table~\ref{tab:ablation}, baseline vs.\ +sql2skeleton row), because
the skeleton abstracts over database-specific column names that appear in
the question but may not match the training database.
The largest structural benefit accrues to \emph{Hard} queries, which gain
+17.1\,pp EX upon adding sql2skeleton
(Section~\ref{sec:exp_difficulty}).


\subsection{LLM Prompt Construction}

The prompt follows the DAIL-SQL format~\citep{gao2023dailsql}: it opens with
\texttt{CREATE TABLE} DDL for the top-$k$ retrieved tables, each accompanied
by three representative sample rows; this is followed by $k$ few-shot
question--SQL pairs ordered by decreasing sql2skeleton similarity; and it
closes with the target question and the \texttt{SELECT} trigger token.
The local Qwen3.5 9B model~\citep{qwen2025} generates SQL completions via
Ollama at \texttt{localhost:11434}.
Total prompt length is approximately 1{,}400 tokens for $k$\,=\,5.

\subsection{Self-Consistency and Two-Pass Correction}

Self-consistency~\citep{wang2022selfconsistency} generates $N$ independent
SQL candidates at temperature $T\,=\,0.7$ and selects the first candidate
that executes without error against the database.
If no candidate executes successfully, the first candidate is returned.
At $N\,=\,3$, self-consistency reduces execution errors by
57.6\% relative to single-pass generation
(Table~\ref{tab:ablation}).
When the selected candidate raises a SQL execution error, the
two-pass corrector re-prompts the LLM with the original prompt, the
failing SQL, and the database error message, reducing errors by a
further 18\%.
SC\,=\,3 addresses structural failures by sampling diverse candidates;
two-pass correction addresses runtime failures (wrong column name,
type mismatch) using error feedback.
